\chapter{Some Calculations of $\operatorname{Gal}(L/K)$}
\section{First Examples}

\begin{proposition} \label{prop-injectToSymmetricGroups}
    Let $L:K$ be the splitting field of a separable polynomial $p(x)\in K[x]$. Let $\alpha_{1},\dots \alpha_{n}\in L$ be the distinct roots of $p(x)$. Then the map $\Phi:\operatorname{Gal}(L /K)\to S_{n}$ defined by $\gamma\mapsto \Phi(\gamma)$ such that $\gamma(\alpha_{i})=\alpha_{\Phi(\gamma)(i)}$ is an injective homomorphism. 
\end{proposition}

\begin{proof}
    $\Phi$ is well-defined since all $\alpha_{i}$ are distinct and it is obviously a group homomorphism. Suppose $\Phi(\gamma)=e_{S_{n}}$, then $\gamma(\alpha_{i})=\alpha_{i}$ for all $i$. But $\gamma:L=K(\alpha_{1},\dots, \alpha_{n})\to L=K(\alpha_{1},\dots, \alpha_{n})$ is uniquely determined by $\gamma(\alpha_{1})=\alpha_{1},\dots, \gamma(\alpha_{n})=\alpha_{n}$. So $\gamma=id_{L}$ and $\ker \Phi = \{id_L\}$. Consequently, $\Phi$ is injective. 
\end{proof}

By abuse of notations, we will treat $\operatorname{Gal}(L /K)$ as a subgroup of $S_{n}$ whenever $L /K$ is the splitting field of some irreducible polynomial of degree $n$ with no repeated roots. 

\begin{example}
    Let $x^{3}-2\in \mathbb{Q}[x]$. Then $\mathbb{Q}(\sqrt[3]{2},\sqrt[3]{2}\omega, \sqrt[3]{2}\omega^{2})=\mathbb{Q}(\sqrt[3]{2},\omega)$ is a splitting field of $x^{3}-2$. 
    So $\mathbb{Q}(\sqrt[3]{2},\omega):\mathbb{Q}$ is Galois with $\left| \operatorname{Gal}(\mathbb{Q}(\sqrt[3]{2},\omega) /\mathbb{Q}) \right|=[\mathbb{Q}(\sqrt[3]{2},\omega):\mathbb{Q}]=6$. By Proposition \ref{prop-injectToSymmetricGroups}, $\operatorname{Gal}(\mathbb{Q}(\sqrt[3]{2},\omega) /\mathbb{Q}) \subseteq S_{3}$. So $\operatorname{Gal}(\mathbb{Q}(\sqrt[3]{2},\omega) /\mathbb{Q}) = S_{3}$. 

    Under $\Phi$ defined in Proposition \ref{prop-injectToSymmetricGroups}, one can write down $\operatorname{Gal}(\mathbb{Q}(\sqrt[3]{2},\omega) /\mathbb{Q}) = S_3$ precisely. Namely,
    let
    \[
    \begin{pmatrix}
        \sqrt[3]{2}\mapsto \sqrt[3]{2}\\
        \omega\mapsto \omega^{2}
    \end{pmatrix} \leftrightarrow (12),\qquad \begin{pmatrix}
        \sqrt[3]{2}\mapsto \sqrt[3]{2}\omega\\
        \omega\mapsto \omega
    \end{pmatrix}\leftrightarrow(123)
    \]
    (note that the elements in $\operatorname{Gal}(\mathbb{Q}(\sqrt[3]{2},\omega) /\mathbb{Q})$ on the left are of order $2$ and $3$ respectively). 
    By Theorem \ref{thm-FTGT-1}, the relation between subextensions and Galois groups are as follows: 
    {\small
    \[
        \begin{tikzcd}[column sep=small]
            && {\mathbb{Q}(\sqrt[3]{2},\omega)} &&&& {\{e\}} \\
            {\mathbb{Q}(\sqrt[3]{2})} & {\mathbb{Q}(\sqrt[3]{2}\omega)} & {\mathbb{Q}(\sqrt[3]{2}\omega^{2})} & {\mathbb{Q}(\omega)} & {\{\langle(12)\rangle\}} & {\{\langle(13)\rangle\}} & {\{\langle(23)\rangle\}} & {\{\langle(123)\rangle\}} \\
            && {\mathbb{Q}} &&&& {S_{3}}
            \arrow[from=1-7, to=2-5]
            \arrow[from=1-7, to=2-6]
            \arrow[from=1-7, to=2-7]
            \arrow[from=1-7, to=2-8]
            \arrow[from=2-1, to=1-3]
            \arrow[from=2-2, to=1-3]
            \arrow[from=2-3, to=1-3]
            \arrow[from=2-4, to=1-3]
            \arrow[squiggly, tail reversed, from=2-4, to=2-5]
            \arrow[from=2-5, to=3-7]
            \arrow[from=2-6, to=3-7]
            \arrow[from=2-7, to=3-7]
            \arrow[from=2-8, to=3-7]
            \arrow[from=3-3, to=2-1]
            \arrow[from=3-3, to=2-2]
            \arrow[from=3-3, to=2-3]
            \arrow[from=3-3, to=2-4]
        \end{tikzcd}
    \]}
    Moreover, for $M=\mathbb{Q}(\sqrt[3]{2})$, $\operatorname{Gal}(L /M)=\{\gamma\in \operatorname{Gal}(L /\mathbb{Q})|\gamma(\sqrt[3]{2})=\sqrt[3]{2}\}=\{e,(12)\}$. $\left| \operatorname{Gal}(L /M) \right| =[L:M]=2$. 

    For $H=\langle(123)\rangle\leq \operatorname{Gal}(L /\mathbb{Q})$, one has $L^{H}=\mathbb{Q}(\omega)$. 

    \medskip
    To verify Theorem \ref{thm-FTGT-2}, note that 
    \begin{center}
        $H=\langle(123)\rangle\lhd S_{3}$ is normal $\iff$ $\mathbb{Q}(\omega) /\mathbb{Q}$ is Galois,
    \end{center} 
    \begin{center}
    $\langle(12)\rangle\leq S_{3}$ is not normal $\iff$ $\mathbb{Q}(\sqrt[3]{2}) :\mathbb{Q}$ is not Galois.     
    \end{center}
\end{example}

\section{Transitivity of $\operatorname{Gal}(L/K)$}

\begin{definition}
    Let $H\leq S_{n}$, we say $H$ is {\bf transitive} if for any $i,j\in \{1,\dots,n\}$, $\exists \sigma\in H$ such that $\sigma(i)=j$. 
\end{definition}

In other words, if we treat $H$ as an $H$-action on $S=\{1,\dots, n\}$, $H$ is transitive means $\mathrm{orb}(i)=S$ for all $i$. 

\begin{example}
    Let $S_{4}$ act on $\{1,2,3,4\}$, then $\langle(1234)\rangle$ is transitive but $\langle(12)\rangle$, $\langle(12)(34)\rangle$ are not. 
\end{example}

\begin{theorem}\label{thm-transitivityGaloisGroup}
    Let $L :K$ be a splitting extension of $p(x)\in K[x]$ with $p(x)$ separable and irreducible. Then $\operatorname{Gal}(L /K)\leq S_{n}$ ($n=\deg p(x)$) is transitive. 
\end{theorem}

\begin{lemma}\label{lem-orbitMinimalPoly}
    Let $L :K$ be Galois, $\alpha\in L$. Then $m_{\alpha}(x)\in K[x]$ satisfies 
    $$m_{\alpha}(x)=\prod_{\beta\in \mathrm{orb}(\alpha)}(x-\beta),$$ where $\mathrm{orb}(\alpha)=\{\gamma(\alpha)|\gamma\in \operatorname{Gal}(L /K)\}$. 
\end{lemma}

\begin{proof}
    Let $\mu(x)=\prod_{\beta\in \mathrm{orb}(\alpha)}(x-\beta)\in L[x]$. Then $\mu(x)\in L^{\operatorname{Gal}(L/K)}[x]=K[x]$ by the proof of $(2) \Rightarrow (3)$ in the proof of Theorem \ref{thm-galcondition}. Moreover, since $\alpha \in \mathrm{orb}(\alpha)$, one has $\mu(\alpha)=0$. So 
    $$m_{\alpha}(x)\mid \mu(x).$$ 
    
    We next prove $\mu(x)$ is irreducible. If not, write $\mu(x)=a(x)\cdot b(x)$ with $a(x), b(x)\in K[x]$ and $\deg a, \deg b<\deg \mu$. Since $\mu(\alpha)=0$, after interchanging $a$ and $b$ if necessary, assume $a(\alpha)=0$, then for any $\beta=\gamma(\alpha)\in \mathrm{orb}(\alpha)$, 
    \[a(\beta)=a(\gamma(\alpha))=\gamma(a(\alpha))=0.\]
    So $\beta$ is also a root of $a(x)$, and hence $(x-\beta) | a(x)$ for all $\beta \in \mathrm{orb}(\alpha)$.
    However $\deg a<\deg \mu=\left| \mathrm{orb}(\alpha) \right| $, contradiction! 
    We must have $\mu(x)$ irreducible and thus 
    $$m_{\alpha}(x) = \mu(x)$$ 
    is the minimal polynomial. 
\end{proof}

\begin{proof}[Proof of Theorem \ref{thm-transitivityGaloisGroup}]
    Replacing $p(x)$ by its monic associate, we may assume that $p(x)$ is monic. Let $\alpha_{1},\dots, \alpha_{n}$ be all the roots of $p(x)$, so that
    $L=K(\alpha_{1},\dots, \alpha_{n}):K$ and $p(x)=\prod_{i=1}^{n}(x-\alpha_{i})\in L[x]$. 
    Since $p(x)$ is irreducible in $K[x]$, $p(x)$ must be the minimal polynomial of each $\alpha_{i}$, for any $1\leq i\leq n$. Therefore, 
    $$\prod_{i=1}^{n}(x-\alpha_{i}) = p(x) = m_{\alpha_i}(x) = \prod_{\beta\in \mathrm{orb}(\alpha_i)}(x-\beta)\in L[x]$$
    for all $1 \leq i \leq n$ by Lemma \ref{lem-orbitMinimalPoly}. By comparing the linear terms on the above equation, one concludes that $\mathrm{orb}(\alpha_{1})=\mathrm{orb}(\alpha_{2})=\cdots=\mathrm{orb}(\alpha_{n}) = \{\alpha_1, \dots, \alpha_n\}$, i.e. $\operatorname{Gal}(L/K)$ is transitive.
\end{proof}

\section{Parity of $\operatorname{Gal}(L/K)$}
\begin{theorem} \label{thm-discrim}
Let $ p(x) \in K[x] $ and $ L/K $ as above. Consider 
$$
\Delta = \Delta(\alpha_1,\ldots,\alpha_n) := \prod_{i<j} (\alpha_i - \alpha_j)^2,
$$ 
where $ \alpha_1,\ldots,\alpha_n $ are roots of $ p(x) $. Then:
\begin{itemize}
    \item $ \Delta \in K^* $, and it can be expressed in terms of coefficients of $ p(x) $
    \item If $ \operatorname{char}(K) \neq 2 $, then
    $$
    \operatorname{Gal}(L/K) \leq A_n \iff \Delta \text{ is a square in } K^*,\ \text{i.e. }\Delta = a^2 \text{ for some } a \in K^*
    $$
\end{itemize}
\end{theorem}

\begin{example}
For $ p(x) = x^3 - 2 \in \mathbb{Q}[x] $, we have $ \Delta = -108 \notin \mathbb{Q}^{*2} $, hence
$$
\operatorname{Gal}(\mathbb{Q}(\sqrt[3]{2}, \omega)/\mathbb{Q}) \nleq A_3.
$$
In fact, $ \operatorname{Gal}(\mathbb{Q}(\sqrt[3]{2}, \omega)/\mathbb{Q}) = S_3 $.
\end{example}

\begin{proof}
    \textbf{(1)} Note that $\Delta$ is symmetric in $\alpha_1,\cdots,\alpha_n$. Since
$\operatorname{Aut}_K(L)\leq S_n$ by Proposition \ref{prop-injectToSymmetricGroups}, 
    $$\Delta\in L^{\operatorname{Aut}_K(L)} = K.$$ 
    By the Fundamental Theorem of Symmetric Polynomials over a field $K$, \[\Delta \in K[s_1(\alpha_1, \alpha_2, ..., \alpha_n), ..., s_n(\alpha_1, \alpha_2, ..., \alpha_n)]\] 
    On the other hand, if $p(x)=c\prod_{i=1}^{n}(x-\alpha_i)$, then
    \[
    \frac{p(x)}{c}=x^n-s_1(\alpha_1,\dots,\alpha_n)x^{n-1}+\cdots+(-1)^ns_n(\alpha_1,\dots,\alpha_n),
    \]
    so every $s_i(\alpha_1,\dots,\alpha_n)$ is determined by the coefficients of $p(x)$, and hence
    $\Delta$ can be expressed by the coefficients of $p(x)$. 
    
\textbf{(2)} Consider
$$
\sqrt{\Delta} := \prod_{i<j} (\alpha_i - \alpha_j) \in L
$$
For all $\sigma \in \operatorname{Gal}(L/K)$, $\sigma$ permutes the $\alpha_i$'s. So one has:
$$
\sigma(\sqrt{\Delta}) = \operatorname{sgn}(\sigma)\sqrt{\Delta}
$$
where $\operatorname{sgn}(\sigma)$ is the permutation signature (here we treat $\sigma \in S_n$). Therefore, when $\operatorname{char}(K) \neq 2$ (so that $1 \neq -1$ in $K$), one has:
\begin{align*}
    \sqrt{\Delta} \in K^* &\iff \sigma(\sqrt{\Delta}) = \sqrt{\Delta}\ \forall \sigma \in  \operatorname{Gal}(L/K)\\ 
    &\iff \operatorname{sgn}(\sigma) = 1\ \forall \sigma \in \operatorname{Gal}(L/K)\\ 
    &\iff \operatorname{Gal}(L/K) \leq A_n
\end{align*}
So the result follows.
\end{proof}

\section{Cyclotomic Extension}  
In this section, we'll study $ x^n - 1 \in K[x] $. Assume $\operatorname{char}(K) = 0$ throughout (ensuring separability).

\begin{definition}
Let $ K $ be any field and $ n \in \mathbb{N} $. The {\bf $n^{th}$ roots of unity in $ K $} are:
$$
\mu_n(K) := \{ \xi \in K \mid \xi^n = 1 \}
$$
\end{definition}
\begin{example}
    \begin{itemize}
        \item $\mu_n(\mathbb{C})=\{e^{2\pi i a/n},0\leq a\leq n-1\}$
        \item $\mu_8(\mathbb{Q}(i))=\{-1,+1,-i,+i\}$
    \end{itemize}
\end{example}

\begin{proposition}
    $\mu_n(K)$ is a cyclic group.
\end{proposition}
\begin{proof}
Note that $\mu_n(K)$ are roots of $x^n - 1$, so there are at most $n$ elements in $\mu_n(K)$, hence $|\mu_n(K)| \leq n$. Obviously, $\mu_n(K)$ is an abelian group.

To see $\mu_n(K)$ is cyclic, we use the fact from MAT3004: For a finite abelian group $G$:
\begin{center}
    \fbox{$G$ is cyclic if and only if 
    for all $d \mid |G|$, there exists a unique subgroup $H \leq G \text{ with } |H| = d$.}
\end{center}


Suppose $|\mu_n(K)| = k$. For any $\delta \in \mu_n(K)$, $\delta^k = 1$, hence:
$$
x^k - 1 = \prod_{\delta \in \mu_n(K)} (x - \delta)
$$
For any $d \mid k$, $x^k-1 \in K[x]$ can be factorized into:
$$
x^k - 1 = (x^d - 1)\left(x^{k-d} + x^{k-2d} + \cdots + 1\right)=\left(\prod_{i=1}^{d}(x-\delta_i)\right)\left(\text{other terms}\right)
$$
Thus the roots $H :=\{\delta\in K\ |\  \delta^d=1 \} = \{\delta_1, \delta_2, \ldots, \delta_d\}$ form a subgroup of $\mu_n(K)$ of order $d$. Moreover, for any $H' \leq \mu_n(K)$ with $|H'| =d$, the elements $\epsilon \in H'$ must satisfy $\epsilon^d = 1$ by Lagrange's Theorem. Hence $\epsilon \in H$ and $H' = H$ is unique.\qedhere
\end{proof}

\begin{remark}
By the above proposition, one has 
\begin{center}
$ x^n - 1 \in K[x]$ splits over $K$ \quad $\Leftrightarrow$ \quad $|\mu_n(K)| = n$ \quad $\Leftrightarrow$ \quad $\mu_n(K) \cong \mathbb{Z}/n\mathbb{Z}$.    
\end{center}
For instance, since $ \mathbb{Q}(i) $ does \emph{not} split $ x^8 - 1 $, so
$
\mu_8(\mathbb{Q}(i)) \cong \mathbb{Z}/4\mathbb{Z}$ but not isomorphic to $\mathbb{Z}/8\mathbb{Z}
$.
\end{remark}

\begin{definition}
Let $K$ be a field, and $L:K$ be such that $\mu_n(L) \cong \mathbb{Z}/n\mathbb{Z}$ (for instance one can take the splitting extension $L:K$ of $x^n-1 \in K[x]$). The {\bf primitive $n^{th}$ roots of unity} $\omega \in \mu_n(L)$ are the generators of the cyclic group satisfying $\mu_n(L) = \langle \omega \rangle$.
\end{definition}

\begin{example} \label{eg-primroots}
    The $n^{th}$ primitive roots of $\mu_n(\mathbb{C})$ are $e^{2\pi i a/n}$ for $\gcd(a,n) = 1$.
    More generally, if $\omega \in \mu_n(L)$ is primitive, then all other primitive roots of $\mu_n(L)$ are of the form
    $$\{\omega^a\ |\ \gcd(a,n) = 1\} = \{\omega^a\ |\ a \in \mathbb{Z}_n^*\}.$$
    In other words, the $n^{th}$ primitive roots are parametrized by $\mathbb{Z}_n^*$.
\end{example}

\begin{definition}
    Let $L:K$ be such that $\mu_n(L) \cong \mathbb{Z}/n \mathbb{Z}$. Define the polynomial
    $$\Phi_{n,K}(x) := \prod_{\beta\ \text{primitive root}}(x-\beta) = \prod_{a \in \mathbb{Z}_n^*}(x-\omega^a)$$
for any primitive $n^{th}$ root $\omega \in \mu_n(L)$.
\end{definition}
\begin{example}
Let $K = \mathbb{Q}$. When $ p $ is prime,
$$
\Phi_p(x) = \prod_{d=1}^{p-1} \left(x - e^{2\pi i d/p}\right) = 1 + x + \cdots + x^{p-1}
$$
When $ n = 4 $:
$$
\Phi_4(x) = (x - i)(x + i) = x^2 + 1 
$$
And when $ n = 6 $:
$$
\Phi_6(x) = (x - \omega)(x - \omega^5) = x^2 - x + 1 \quad (\omega = e^{2\pi i/6})
$$
\end{example}
In fact, $\Phi_{n,K}(x)$ is independent of the choice of the splitting field extension $L:K$. This justifies our symbol $\Phi_{n,K}(x)$ (rather than $\Phi_{n,L}(x)$).
\begin{lemma}
    $\Phi_{n,K}(x)\in K[x] $, and is independent of the choice of $L:K$ satisfying $\mu_n(L) \cong \mathbb{Z}/n\mathbb{Z}$.
\end{lemma}
\begin{proof}
    For any such $L:K$, let $\omega \in \mu_n(L)$ be a primitive root of unity. Then we have a splitting field $K\subseteq K(\omega) \subseteq L$ where $K(\omega)$ is the splitting field of $x^n-1 = \prod_{i \in \mathbb{Z}_n}(x-\omega^i)$. Therefore $K(\omega):K$ is Galois, and 
    $$\Phi_{n,K}(x) = \prod_{a \in \mathbb{Z}_n^*}(x-\omega^a) \in K(\omega)[x].$$ 
    Note that $\forall \sigma \in \operatorname{Gal}(K(\omega)/K)$, $\sigma$ maps primitive root to primitive root, that is $\sigma$ permutes the primitive roots of unity. 
    Since the coefficients of $\Phi_{n,K}(x)$ are symmetric over the primitive roots, they are all fixed by $\sigma \in \operatorname{Gal}(K(\omega)/K)$. In other words, the coefficients of $\Phi_{n,K}$ are 
    in $K(\omega)^{\operatorname{Gal}(K(\omega)/K)}=K$. If a different splitting field is chosen, uniqueness of splitting fields gives a $K$-isomorphism carrying primitive roots to primitive roots, so the same polynomial is obtained.
\end{proof}

Due to the importance of $K(\omega):K$, we give the following:
\begin{definition}
The splitting extension $K(\omega):K$ of $x^n-1 \in K[x]$ is called the {\bf $n^{th}$ cyclotomic extension} of $K$.   
\end{definition}

\begin{theorem}
Let $ K(\omega)/K $ be $n^{th}$ cyclotomic extension. Then:
\begin{enumerate}
    \item There exists an injective homomorphism
    $$
    \phi: \operatorname{Gal}(K(\omega)/K) \hookrightarrow \mathbb{Z}_n^\star
    $$
    
    \item $\phi$ is an isomorphism if and only if the cyclotomic polynomial $\Phi_{n,K}(x)$ is irreducible over $K$.
\end{enumerate}
\end{theorem}
\begin{proof}
    \begin{enumerate}
        \item Since $\phi$ maps primitive roots to primitive roots, then $\sigma(\omega)=\omega^{a(\sigma)}$, where $a(\sigma)\in \mathbb{Z}_n^*$, we define: $$\phi:\operatorname{Gal}(K(\omega)/K) \rightarrow \mathbb{Z}_n^\star \quad \text{by} \quad \phi(\sigma):=a(\sigma).$$
        It is easy to check that this is an injective homomorphism.
        \item $(\Rightarrow)$ Suppose $\phi$ is an isomorphism, then $ \operatorname{Gal}(K(\omega)/K) \simeq \mathbb{Z}_n^\star $, thus $\operatorname{orb}(\omega)=\{\omega^{a} \ |\ a \in \mathbb{Z}_n^\star \}$ consists all primitive roots (Example \ref{eg-primroots}). So
        $$\Phi_{n,K}(x)=\prod_{\beta \text{ primitive roots}} (x-\beta) = \prod_{\omega^a \in \operatorname{orb(\omega)}}(x-\omega^a)=m_\omega(x),$$
        where the last equality comes from Lemma \ref{lem-orbitMinimalPoly}. But we know $m_\omega  (x) $ is irreducible, so $\Phi_{n,K}(x)$ is also irreducible.
        \\
        $(\Leftarrow)$ Suppose $\Phi_{n,K}(x) $ is irreducible, then as above, we have 
        $$\Phi_{n,K}(x)=m_\omega(x)=\prod_{\beta \text{ primitive  roots}}(x-\beta) =  \prod_{\omega' \in \operatorname{orb(\omega)}}(x-\omega')$$
        then $|\operatorname{orb}(\omega)|=|\mathbb{Z}_n^\star|$. Since an automorphism of $K(\omega)$ fixing $\omega$ is the identity, the stabilizer of $\omega$ is trivial. Hence
        \[
        |\operatorname{Gal}(K(\omega)/K)|=|\operatorname{orb}(\omega)|=|\mathbb{Z}_n^\star|.
        \]
        Since $\phi$ is injective, it is an isomorphism.
    \end{enumerate}
\end{proof}

For instance, if $K = \mathbb{C}$, then any $\Phi_{n,\mathbb{C}}(x)$ splits in $\mathbb{C}[x]$, i.e. $\Phi_{n,\mathbb{C}}(x)$ is reducible if its degree is $> 1$, and hence $\phi: \operatorname{Gal}(\mathbb{C}(\omega)/\mathbb{C}) \to \mathbb{Z}_n^*$ is almost never an isomorphism (which is obviously true since $[\mathbb{C}(\omega):\mathbb{C}] = 1$. On the other extreme, if $K = \mathbb{Q}$, one has:
\begin{theorem}
    $\Phi_{n,\mathbb{Q}}(x)\in \mathbb{Z}[x]$ is always irreducible. Consequently, $\operatorname{Gal}(\mathbb{Q}(\omega)/\mathbb{Q}) \cong \mathbb{Z}_n^\star$.
\end{theorem}
\begin{proof}
    We need to show that $m_\omega(x)=\Phi_{n,\mathbb{Q}}(x)$. Note that $m_\omega(x)|x^n-1$ in $\mathbb{Q}[x]$, so we have 
    $$x^n-1=m_\omega(x) \cdot t(x)$$
    for some $t(x) \in \mathbb{Q}[x]$.
    \begin{itemize}
        \item Claim 1: $m_\omega(x), t(x) \in \mathbb{Z}[x]$
        \item Claim 2: For any prime $p$ such that $\operatorname{gcd}(p,n)=1$, $m_{\omega}(\omega^p)=0$
    \end{itemize}
    Suppose the claims hold, then for any $a$ such that $\operatorname{gcd}(a,n)=1$, $a=p_1p_2\cdots p_n$ with $\gcd(p_i,n) = 1$. Then 
    \begin{align*}
    m_\omega(\omega)=0 &\Rightarrow m_\omega(\omega^{p_1})=0\ (\text{so } m_{\omega} = m_{\omega^{p_1}}) \\ &\Rightarrow m_\omega(\omega^{p_1p_2})=m_{\omega^{p_1}}((\omega^{p_1})^{p_2})=0 \\ 
    &\ \ \vdots \\
    &\Rightarrow m_\omega(\omega^{p_1p_2\dots p_n}) = m_\omega(\omega^a)=0\end{align*}
    and then we conclude that $m_\omega(x) $ contains all the primitive roots $\omega^a$ for $\gcd(a,n) =1$, i.e.
    $$\Phi_{n,\mathbb{Q}}(x)\ |\ m_\omega(x)$$
    and the equality comes from the fact that $m_{\omega}(x)$ is irreducible.
    
    \medskip
    We left Claim 1 as an exercise (it uses the proof of Gauss Lemma), now we prove Claim 2 -- Suppose on contrary that $m_\omega(\omega^p)\neq 0$. Since
    \[
    m_\omega(\omega^p)t(\omega^p)=(\omega^p)^n-1=0,
    \]
    we have $t(\omega^p)=0$, and hence $m_\omega(x)\mid t(x^p)$ in $\mathbb{Z}[x]$. Thus there exists $h(x)\in\mathbb{Z}[x]$ such that
    \[
    m_\omega(x)h(x)=t(x^p).
    \]
    Reducing modulo $p$ gives
    \[
    \overline{m_\omega}(x)\,\overline{h}(x)=\overline{t}(x)^p
    \qquad\text{in }\mathbb{F}_p[x].
    \]
    Hence $\overline{m_\omega}(x)$ and $\overline{t}(x)$ have a nonconstant common factor $g(x)$. Since
    \[
    x^n-1=\overline{m_\omega}(x)\,\overline{t}(x)
    \qquad\text{in }\mathbb{F}_p[x],
    \]
    we get $g(x)^2\mid x^n-1$. But $\gcd(p,n)=1$, so $x^n-1$ has no repeated roots in $\mathbb{F}_p[x]$, a contradiction.
\end{proof}
\section{Kummer Extension}
In this section, we will study the polynomial $p(x)=x^n-a \in L[x]$, where $L$ satisfies 
\begin{itemize}
    \item $\operatorname{char}(L) \nmid n$ (to ensure $p(x)$ has no repeated roots); and
    \item $x^n-1$ splits in $L[x]$ (we will write $\omega \in L$ as an $n^{th}$ primitive root).
\end{itemize}
Assume $a\in L^\times$. Suppose $M:L$ is the splitting field of $x^n-a$, and let $\beta\in M$ be a root. Then all its other roots are given by $\omega^i\beta \in M$.
\begin{lemma} \label{lem-kummer}
    Let $M:L$ be the splitting field of $x^n-a\in L[x]$. For any root $\rho \in M$ of $x^n-a$, the map 
    $$\psi: \operatorname{Gal}(M/L) \hookrightarrow \mu_n(L)\simeq \mathbb{Z}_n$$ 
$$\psi(\sigma):=\sigma(\rho)/\rho  $$  
is a well-defined injective group homomorphism.

Moreover, if $x^n-a$ is irreducible, then $\psi$ is an isomorphism with $\operatorname{Gal}(M/L) \cong \mathbb{Z}_n$.
\end{lemma}
\begin{proof}
Firstly, we need to show \(\sigma(\rho)/ \rho \in \mu_n(L)\). Since \begin{align*}
    \left(\frac{\sigma(\rho)}{\rho} \right)^n = \frac{\sigma(\rho^n)}{\rho^n} = \frac{\sigma(a)}{a} = \frac{a}{a} = 1,
\end{align*}  
the last two equalities holds since \(\sigma \in \mathrm{Gal}(M/L), a \in L\). We observe that \(\sigma(\rho)/\rho \in M\) is a root of the polynomial \(x^n -1\). Since \(x^n -1 = \displaystyle\prod_{1 \leq i \leq n} (x-\omega^i) \in L[x] \subseteq M[x]\).  Therefore \( \sigma(\rho)/\rho = \omega^i \in \mu_n(L) \subseteq \mu_n(M)\).

Secondly, we need to verify our definition of $\psi$ is independent of the choice of \(\rho\) -- Suppose that \(\chi \in M\) also satisfies \(\chi^n = a\), then \(\rho/ \chi = \omega^i \in \mu_n(L)\). Since \((\rho/\chi)^n = a/a =1\), we can compute
\begin{align*}
    \frac{\sigma(\rho)}{\rho} = \frac{\sigma(\omega^i \chi)}{\omega^i \chi} = \frac{\sigma(\omega^i)\sigma(\chi)}{\omega^i\chi} = \frac{\omega^i\sigma(\chi)}{\omega^i\chi} = \frac{\sigma(\chi)}{\chi},
\end{align*}
with \(\theta' \in L\), the \(\sigma(\rho)/\rho\) is well defined.

Thirdly, we need to check that \(\psi\) is a group homomorphism: For all \(\sigma_1 , \sigma_2 \in \mathrm{Gal}(M/L)\),
\begin{align*}
     & \quad \quad \quad \psi(\sigma_1 \sigma_2) = \psi(\sigma_1) \psi(\sigma_2)\\
     &\Leftrightarrow \quad\frac{\sigma_1\sigma_2(\rho)}{\rho} = \frac{\sigma_1(\rho)}{\rho} \cdot\frac{\sigma_2(\rho)}{\rho}\\
    &\Leftrightarrow \quad  \sigma_1 \sigma_2(\rho) = \frac{\sigma_1(\rho) \sigma_2(\rho)}{\rho}\\
     &\Leftrightarrow \quad \frac{\sigma_1 \sigma_2(\rho)}{\sigma_1(\rho)} = \frac{\sigma_2(\rho)}{\rho}\\
   &\Leftrightarrow \quad 
\sigma_1\left(\frac{\sigma_2(\rho)}{\rho}\right) = \frac{\sigma_2(\rho)}{\rho}.
\end{align*}
Note that \(\sigma_2(\rho)/\rho \in \mu_n(L)\), so \(\sigma_1\in \mathrm{Gal}(M/L)\) fixes \(\sigma_2(\rho)/\rho\) and hence the last equality holds.

Finally, to show \(\psi\) is injective, suppose that \(\psi(\sigma) = 1 \in \mu_n(L)\), then \(\sigma(\rho)/\rho = 1\) and hence \(\sigma(\rho) = \rho\). For all roots \(\left\{ \rho \omega^i| 0 \leq i \leq n-1\right\}\) of polynomial \(x^n -a\), \(\sigma(\rho \omega^i) = \sigma(\rho) \sigma(\omega^i) = \rho \omega^i\), since \(\omega^i \in L\). Therefore, \(\sigma\) fixes all roots of \(x^n-a\), so it is the identity map on \(M=L(\rho,\rho\omega,\dots,\rho\omega^{n-1})\).

If $x^n-a$ is irreducible, then $M=L(\rho)$ and $[M:L]=n$. Also $|\mu_n(L)|=n$. Hence the injective homomorphism $\psi$ is an isomorphism.
\end{proof}

The importance of the Lemma \ref{lem-kummer} is that its converse also holds:
\begin{theorem}\label{thm-cyclicGalImpliesKummerExtension}
    Let $L$ be such that $x^n-1$ splits in $L$ and $\operatorname{char}(L)\nmid n$. Suppose  $M:L$ is a Galois extension such that $\operatorname{Gal}(M/L) = \langle \sigma \rangle \cong \mathbb{Z}_n$. Let $\omega \in L$ be a primitive $n^{th}$ root of unity, then
\begin{enumerate}
    \item[(a)] There exists $\alpha \in M$ such that
    $$\beta := \alpha + \omega\sigma(\alpha) + \dots + \omega^{n-1}\sigma^{n-1}(\alpha) \neq 0.$$
    \item[(b)] $b := \beta^n \in L$.
    \item[(c)] $M = L(\beta)$ is a simple extension, with $m_{\beta}(x) = x^n - b$.
\end{enumerate}
\end{theorem}
From now on, we call such $M:L$ satisfying the hypotheses of the above theorem a {\bf Kummer extension}.
\begin{proof}
    (a) is proved in Homework 2 Question 8. For (b), one can easily check that \(\sigma(\beta) = \omega^{-1} \beta \) which implies
    \begin{equation*}
        \sigma^i(\beta) = \omega^{-i} \beta     \quad \quad
    \end{equation*}
    Also, \(\sigma^i(\beta^n) = \left(\sigma^i(\beta)\right)^n = \omega^{-in} \beta^n = 1^{-i}\beta^n = \beta^n\). Therefore, \(\beta^n\) is fixed by \(\langle \sigma \rangle = \mathrm{Gal}(M/L)\), i.e. \(\beta^n \in M^{\mathrm{Gal}(M/L)} = L.\)

    For (c), Obviously, \(L(\beta) \subseteq M\). Since \([M:L] = |\mathrm{Gal}(M/L)| = n\), by tower law, we just need to show \(\left[L(\beta):L\right]=n\). We study the minimal polynomial \(m_\beta(x) \in L[x]\), by lemma \ref{lem-orbitMinimalPoly}
    \begin{align*}
        m_\beta(x) = \prod_{\gamma \in \mathrm{orb}(\beta)} (x- \gamma).
    \end{align*}
    But Equation $(*)$ implies that \(\mathrm{orb}(\beta) = \left\{\sigma^i(\beta)\right\} = \{\beta, \omega^{-1}\beta, \dots, \omega^{-(n-1)}\beta\}\), then
    \begin{align*}
        m_{\beta}(x) =& \prod_{i=0}^{n-1}\left(x - \sigma^{i}(\beta)\right) = \prod_{i=0}^{n-1}(x-\omega^{-i}\beta) =x^n - \beta^n = x^n - b
    \end{align*}
    Therefore, \([L(\beta):L] = \mathrm{deg}\left(m_\beta(x)\right) = n\).
\end{proof}
Before Theorem \ref{thm-cyclicGalImpliesKummerExtension}, we compute the Galois group $\operatorname{Gal}(M/L)$ of a given field extension $M:L$. However, Theorem \ref{thm-cyclicGalImpliesKummerExtension} works in the opposite direction -- we obtain the field extension $M:L$ from its Galois group $\operatorname{Gal}(M/L)$. This will play an important role in the study of roots of polynomials in the next section.
